\chapter{Background}
\label{chap:background}

This chapter provides the background required to understand the rest of the
thesis. It is not intended to be a complete introduction to compiler theory.
Instead, it explains the concepts that directly influenced the design of \epl:
ahead-of-time compilation, lexical analysis, parsing, static typing, intermediate
representations, control-flow graphs, SSA-like values, compile-time evaluation,
foreign function interfaces, and LLVM.

\section{Compiled Languages}

A compiler translates a program from one language into another language. In a
systems programming context, the source language is usually a human-readable
programming language and the output is usually machine code, assembly, object
code, bytecode, or another lower-level representation. Classic compiler texts
describe this as a sequence of front-end, middle-end, and back-end phases
\cite{aho2006compilers,cooper2011engineering}. The front end understands source
syntax and semantics. The middle end performs analysis and transformations. The
back end emits code for the target execution environment.

An ahead-of-time compiler performs this translation before the program is run.
This differs from an interpreter, which executes a program directly, and from a
just-in-time compiler, which compiles parts of the program while it is already
running. Ahead-of-time compilation is suitable for \epl{} because the project
targets native object output and C-compatible linking. The compiler performs all
source analysis before execution, emits an object file, and leaves final linking
to a system C compiler.

Ahead-of-time compilation has several advantages for a systems language. Errors
can be detected early, generated code can be inspected, external symbols can be
linked by existing toolchains, and the final executable does not need to include
an interpreter. It also has costs. Compilation must know enough about all
referenced functions and types up front, compile-time execution must be
restricted to information available at compilation, and the implementation must
make target-specific decisions about data layout and calling conventions.

\section{Lexical Analysis}

Lexical analysis is the first stage of most compilers. The lexer reads raw source
text and produces a sequence of tokens. Tokens are the meaningful units of the
surface language: identifiers, keywords, literals, operators, delimiters, and
punctuation. Whitespace and comments are usually removed at this stage unless
they are semantically important.

The \epl{} lexer is hand-written. It recognizes identifiers, reserved keywords,
number literals, string literals, punctuation, whitespace, and comments beginning
with the \texttt{\#} character. It also attaches spans to tokens. A span is the
start and end byte offset of a source fragment. Later stages use spans to report
diagnostics at the correct line and with a caret underline.

Using a hand-written lexer has an educational advantage: the exact rules are
visible in ordinary Rust code. It also matches the scale of the language. \epl{}
does not require indentation-sensitive lexing, contextual lexing, Unicode
identifier normalization, or complex preprocessor behavior. A compact iterator
over tokens is sufficient.

\section{Parsing and Abstract Syntax Trees}

Parsing turns tokens into a tree that reflects the grammar of the language. The
abstract syntax tree, or AST, records the structure of a source program without
yet resolving every semantic question. For example, in the AST a type name is
still an identifier string. A variable use is still a name. A binary expression
has a left operand, an operator, and a right operand, but it does not yet know
whether the operands are integers or booleans.

\epl{} uses recursive descent parsing. In recursive descent, grammar rules are
implemented as functions that call one another. This style is straightforward for
languages whose grammar can be parsed with a small amount of lookahead. It is
also easy to connect expression precedence to function structure: assignment
calls logical-or parsing, logical-or calls logical-and parsing, comparison calls
range parsing, range calls addition, addition calls multiplication, and so on.
The approach is common in educational compilers because it keeps parsing logic
close to the grammar \cite{nystrom2021crafting}.

The \epl{} AST deliberately stays close to source syntax. It stores annotations,
function declarations, struct definitions, block expressions, array and struct
initializers, control-flow expressions, casts, calls, assignments, and all
source-level operators. Most semantic checks are delayed until typed lowering.
This separation keeps parsing focused on syntax and keeps type checking focused
on meaning.

\section{Static Typing}

A statically typed language assigns a type to each expression before the program
is executed. Static typing helps catch errors early, documents programmer intent,
and provides information needed for code generation. In a compiler targeting
native code, type information is also needed to choose instruction widths,
compute memory layouts, decide calling conventions, and validate operations.

\epl{} is statically typed and intentionally avoids implicit numeric
conversions. Signed and unsigned integer types of the same width are distinct in
the typed tree. This decision prevents surprising conversions and makes type
checking simple: arithmetic operands must have the same integer type, casts must
be explicit, and a function call argument must match the parameter type. The
lower IR later collapses signed and unsigned integers of the same width into the
same storage type while preserving signedness on operations that need it.

The language also includes two special types. The \texttt{unit} type is a
zero-sized type used for expressions with no meaningful value, similar to
\texttt{void} in C but represented as an ordinary value type in the compiler. The
never type, written \texttt{!}, is used for expressions that do not complete
normally, such as \texttt{return}, \texttt{break}, and calls to external
functions such as \texttt{exit}. Treating never as a type simplifies the typing
of expression-oriented control flow.

\section{Intermediate Representations}

An intermediate representation, or IR, is a compiler-internal representation of
a program. Compilers usually use more than one IR because different stages need
different shapes of information \cite{muchnick1997advanced}. A source-like tree
is convenient for type checking and diagnostics, while a control-flow graph is
more suitable for code generation and optimization.

\epl{} uses two main intermediate representations. The first is \irtree{}, a
typed expression tree. It is source-like enough to be readable but semantic
enough that every expression has a type, variables and functions are represented
by generated IDs, aggregate layouts have been computed, and several source
constructs have been lowered into simpler forms. The second is lower IR, a
control-flow graph with basic blocks, instructions, terminators, definitions,
and block arguments. Lower IR is closer to LLVM and makes branches, returns,
loads, stores, calls, pointer offsets, casts, and arithmetic explicit.
Table~\ref{tab:ir-summary} summarises the four representations and their roles.

\begin{table}[H]
  \centering
  \begin{tabular}{p{0.23\textwidth}p{0.30\textwidth}p{0.32\textwidth}}
    \toprule
    Representation & Main purpose & Important properties \\
    \midrule
    AST & Preserve source structure after parsing & Names and type names are still unresolved; spans are kept for diagnostics. \\
    \irtree{} & Semantic analysis and compile-time evaluation & Every expression is typed; variables/functions/loops use IDs; source constructs are simplified. \\
    Lower IR & Code generation and CFG cleanup & Basic blocks, instructions, terminators, block arguments, explicit memory operations. \\
    LLVM IR & Target-independent backend representation & Verified by LLVM and emitted as native object code. \\
    \bottomrule
  \end{tabular}
  \caption{Intermediate representations used by the \epl{} compiler}
  \label{tab:ir-summary}
\end{table}

\section{Control-Flow Graphs and SSA-Like Values}

A control-flow graph represents a function as basic blocks connected by
terminators. A basic block is a sequence of instructions with one entry and one
terminator. Terminators include jumps, conditional jumps, returns, and
unreachable markers. CFGs make the possible execution paths through a function
explicit.

Static Single Assignment form, or SSA, represents each computed value as being
assigned exactly once. SSA is widely used because it simplifies many analyses
and optimizations \cite{cytron1991ssa}. \epl{} does not implement a full
source-variable SSA construction pass. Local mutable variables are represented as
stack allocation slots in lower IR, and loads and stores are explicit. However,
computed lower IR values are immutable definitions, and basic block arguments
serve a similar role to phi nodes at join points. This is a common way to model
SSA-like control flow: predecessor blocks pass values to successor block
arguments, and LLVM generation turns those block arguments into LLVM phi nodes.

\section{Compile-Time Evaluation}

Compile-time evaluation means executing some part of a program while compiling
it, then using the result in the generated program. Partial evaluation and
metaprogramming have a long history in programming language research
\cite{jones1993partial}. Modern systems languages use compile-time computation
for constants, generic specialization, code generation, configuration, and data
precomputation.

Compile-time evaluation is attractive because it moves work out of the final
executable. It can also let programmers write ordinary code to compute constants
instead of using external generators. The difficulty is that not all operations
make sense at compile time. Reading a run-time variable, dereferencing a raw
pointer, calling an impure external function, performing I/O, or depending on
target process state would make the compiler's behavior unsafe or
non-deterministic.

\epl{} handles this by defining a pure compile-time subset. The subset is not a
separate language. It is the same expression language after type checking, but
with a purity check before interpretation. This design avoids building a second
parser or evaluator for a separate macro language.

\section{LLVM}

LLVM is a compiler infrastructure project that provides a typed intermediate
representation, optimization passes, target backends, object emission, and tools
\cite{lattner2004llvm,llvm_langref}. Many language implementations use LLVM to
avoid writing a complete machine-code backend. A compiler can lower its own IR to
LLVM IR, verify the LLVM module, and ask LLVM to emit assembly or object code
for a target platform.

\epl{} uses LLVM through the Rust \texttt{llvm-sys} crate. The compiler declares
all functions first, emits each function body into LLVM basic blocks, creates
allocas for local storage, emits loads and stores, maps lower IR block arguments
to LLVM phi nodes, emits integer arithmetic and comparisons, and finally asks
LLVM to emit a native object file. The object file is then linked with a system
C compiler so that external functions such as \texttt{printf} or \texttt{exit}
can be resolved from libc.
